\subsection{\texorpdfstring{特征为 \(p\) 的交换环}{特征为 p 的交换环}}

在本节及下一节中，p 表示一个素数。

定理 2. --- 设 A 是一个特征为 p 的交换环。映射 \(a \mapsto a^{p}\) 是环 A 的一个自同态，即我们有关系式

\[
(a + b) ^ {p} = a ^ {p} + b ^ {p}\tag{1}
\]

\[
(a b) ^ {p} = a ^ {p} b ^ {p}\tag{2}
\]

对 A 中的 \(a, b\) 成立。

公式 (2) 由 A 的交换性得出。为证明 (1)，我们使用二项式公式 \((a + b)^p = a^p + b^p + \sum_{i=1}^{p-1} \binom{p}{i} \cdot a^ib^{p-i}\)；由于 \(p \cdot x = 0\) 对所有 \(x \in \mathbf{A}\) 成立，只需证明以下引理：

引理 1. --- 设 \(p\) 是一个素数，\(i\) 是从 1 到 \(p - 1\) 范围内的一个整数，则二项式系数 \(\binom{p}{i}\) 是一个可被 \(p\) 整除的整数。

我们对 \(i\) 作归纳论证，情形 \(i = 1\) 由公式 \(\binom{p}{1} = p\) 直接得出。假设 \(2 \leqslant i \leqslant p - 1\) 且 \(\binom{p}{i - 1}\) 能被 \(p\) 整除。则整数 \(i\binom{p}{i} = (p - i + 1)\binom{p}{i - 1}\) 属于素理想 \(p\mathbf{Z}\)（\(\mathbf{Z}\) 的素理想）；由于 \(i \notin p\mathbf{Z}\)，我们有 \(\binom{p}{i} \in p\mathbf{Z}\)，引理得证。

设 A 是特征为 p 的交换环，f 为一个整数 \(\geqslant0\)。由定理 2，我们通过对 f 的归纳法推出，映射 \(a\mapsto a^{p^{f}}\) 是环 A 的一个自同态。特别地，我们有关系式

\[
(a _ {1} + \dots + a _ {n}) ^ {p ^ {f}} = a _ {1} ^ {p ^ {f}} + \dots + a _ {n} ^ {p ^ {f}}\tag{3}
\]

对 A 中任意的 \(a_{1},\ldots,a_{n}\) 成立。映射 \(a\mapsto a^{p}\) 有时称为 A 的弗罗贝尼乌斯自同态。取 \(\mathbf{A}=\mathbf{F}_{p}\) 且 \(a_{i}=1\)，我们由 (3) 得到关系式：

\[
n ^ {p ^ {f}} \equiv n \mod p (n \in \mathbf {Z}, f \in \mathbf {N}).\tag{4}
\]

对 A 的每个子集 S，我们用 \(S^{p^{f}}\) 表示 A 中形如 \(x^{p^{f}}\)（其中 \(x \in S\)）\(^{1}\) 的元素所组成的集合。特别地，若 K 是 A 的一个子环，则集合 \(K^{p^{f}}\) 是 A 的一个子环。若 K 是 A 的一个子环且 S 是 A 的一个子集，我们用 K{[}S{]} 表示由 \(K \cup S\) 生成的 A 的子环；当 A 是一个域时，我们用 \(\mathrm{K}(S)\) 表示 K{[}S{]} 的分式域，即由 \(K \cup S\) 生成的 A 的子域。

命题 2. --- 设 A 是特征为 p 的交换环，K 是 A 的一个子环，S 是 A 的一个子集，f 是一个正整数。

\begin{enumerate}
\def\labelenumi{\alph{enumi})}
\item
  我们有 \(\mathbf{K}[\mathbf{S}]^{p^f} = \mathbf{K}^{p^f}[\mathbf{S}^{p^f}]\)，且若 \(\mathbf{A}\) 是一个域，则 \(\mathbf{K}(\mathbf{S})^{p^f} = \mathbf{K}^{p^f}(\mathbf{S}^{p^f})\) 。
\item
  若 K-模 K{[}S{]} 由 A 的元素组成的族 \((a_{i})_{i\in I}\) 生成，则 K-模 \(K[S^{p^{f}}]\) 由族 \((a_{i}^{p^{f}})_{i\in I}\) 生成。
\end{enumerate}

由于 K{[}S{]} 是 A 中由 K ∪ S 生成的子环，它的像 K{[}S{]} \(^{p^{f}}\) （在环 A 的自同态 \(\pi: a \mapsto a^{p^{f}}\) 下）是 A 中由 K ∪ S 的像 K \(^{p^{f}}\) ∪ S \(^{p^{f}}\) （在 \(\pi\) 下）生成的子环，因此 K{[}S{]} \(^{p^{f}}\) = K \(^{p^{f}}\) {[}S \(^{p^{f}}\) {]}。域的情形类似处理；这就证明了 a)。

显然，族 \((a_i^{p^f})_{i\in I}\) 生成 \(\mathbf{K}^{p^f}\) -模 \(\mathbf{K}[\mathbf{S}]^{p^f}\) 。K-模 \(\mathbf{K}[\mathbf{S}^{p^f}]\) 由形如 \(x_1^{p^f}\ldots x_n^{p^f} = (x_1\ldots x_n)^{p^f}\)（其中 \(x_1,\ldots ,x_n\) 在 S 中任意）的乘积生成，因此也由集合 \(\mathbf{K}[\mathbf{S}]^{p^f}\) 生成。断言 b) 直接由此得出。

\(^{1}\) 当然，集合 \(S^{p^{f}}\) 不应与 \(p^{f}\) 个都等于 S 的集合的集合积相混淆，也不应与属于 S 的 \(p^{f}\) 个元素的乘积组成的集合相混淆。

